% 同一教学规则的四种输入与当前样本；正文宽度设计。
\begin{tikzpicture}[foundation picture]
 \path[use as bounding box] (0,0) rectangle (110,52);
 \node at (25,46) {$f=1+2x_1+4x_2$};
 \node at (7,36) {输入}; \node at (25,36) {分数}; \node at (41,36) {动作};
 \foreach \y/\xone/\xtwo/\score/\result in {28/0/0/1/放行,20/1/0/3/放行,12/0/1/5/放行,4/1/1/7/隔离}{
  \node at (7,\y) {$(\xone,\xtwo)$};
  \node at (25,\y) {$\score$}; \node at (41,\y) {\result};
 }
 \node at (83,46) {当前邮件$(1,1)$};
 \draw[foundation axis,line width=.9pt] (61,9)--(107,9);
 \foreach \x/\v in {61/0,79/3,97/6}{\node[below] at (\x,9) {$\v$};}
 \fill[FigureYellow] (61,25) rectangle (67,32);
 \fill[FigureBlue] (67,25) rectangle (79,32);
 \fill[FigureRed] (79,25) rectangle (103,32);
 \draw[FigureStroke,line width=1pt] (61,25) rectangle (103,32);
 \node[above] at (64,32) {$1$}; \node[above] at (73,32) {$+2$}; \node[above] at (91,32) {$+4$};
 \draw[FigureStroke,dashed,line width=1pt] (97,9)--(97,38);
 \node[anchor=west] at (99,40) {$6$};
 \node[font=\figurefont\fontsize{8}{10}\selectfont] at (64,19) {基准};
 \node[font=\figurefont\fontsize{8}{10}\selectfont] at (75,19) {支付};
 \node[font=\figurefont\fontsize{8}{10}\selectfont] at (91,19) {附件};
\end{tikzpicture}
